\begin{table}[t]\centering
\textbf{[float64 Jacobian run matrix: artifact not present at build time]}
\caption{float64 Jacobian run matrix --- pending.}
\label{tab:jacruns}
\end{table}


\paragraph{Why a degenerate run happens.} At certain sample points a large
fraction of the candidate rows evaluate to numerical zero. Their surviving
entries are rounding noise. Row-normalising rescales that noise to unit length,
after which the singular spectrum shows no gap and the reported rank becomes an
artefact of the tolerance. The rule enforced here is that a row whose norm is at
or below the declared noise floor is never normalised; the withheld
normalisation, and the reason, are recorded in the run record.

\paragraph{Why the exact computation is not subject to this.} The analytic
amputated derivative involves no step size, and rank over $\mathbb{F}_p$ involves
no tolerance. Because the gamma-traceless basis is integral --- entries in
$\{-1,0,+1\}$, annihilated exactly over $\mathbb{Z}$ --- the Jacobian is the
reduction of an integer matrix, so its modular rank is a rigorous lower bound on
the rank in characteristic zero.

\paragraph{Scope.} The exact computation covers the complete archived candidate
selection: $\exactJacScheduled$ scheduled, $\exactJacEvaluated$ evaluated,
$\exactJacErrors$ evaluation errors, $\exactJacZeroRows$ zero rows. Every
candidate carries a terminal status with a value, a derivative row, an output
hash, a runtime and a peak-RSS figure, so a candidate that failed would be
visible rather than silently absent from a rank.

An intermediate version of this appendix quoted $\archiveJacobianRank$ for the
port-graph subset alone ($\nArchivePortGraphs$ of $\nArchiveCandidates$), with
the $\nArchiveStructured$ structured tensor-word candidates excluded rather than
assumed inert. Those candidates were subsequently implemented in the same exact
arithmetic; section~\ref{sec:cert81} reports the result for the full selection.

\paragraph{What the number of primes does and does not buy.} For the rank
statement the prime count is not what carries the argument. An integer matrix
whose determinant vanishes vanishes modulo every prime, so a single prime at
which the explicit $\minorSize\times\minorSize$ minor has non-vanishing
determinant already establishes that the integer minor is nonzero, and hence the
characteristic-zero bound. No height bound, Hadamard estimate or rational
reconstruction enters. Additional primes are run because they guard against an
indexing error in selecting the minor --- a different failure mode, and one a
single prime would not catch. This is in contrast to the subspace dimensions of
section~\ref{sec:decomp}, which are equalities rather than one-sided bounds and
where additional primes do reduce genuine risk.
